\documentclass[11pt]{article}
\usepackage[margin=0.85in]{geometry}
\usepackage{amsmath,amssymb,amsthm}
\usepackage[T1]{fontenc}
\usepackage{lmodern}
\usepackage{xcolor}
\usepackage{hyperref}
\hypersetup{colorlinks=true,linkcolor=blue,urlcolor=blue}
\newtheorem{proposition}{Proposition}
\newcommand{\R}{\mathbb R}
\newcommand{\uu}{\mathbf u_0}
\newcommand{\dd}{\,\mathrm d}
\setlength{\parindent}{0pt}
\setlength{\parskip}{0.55em}
\title{A smooth initial velocity field with swirl}
\author{Working note}
\date{5 October 2026}

\begin{document}
\maketitle

This note completes the initial-data construction: the field below is smooth
through the symmetry axis, divergence-free on all of $\R^3$, and rapidly
decreasing. Its kinetic energy has an exact formula. These statements concern
the starting field. The discrete experiments below do not establish a theorem
about time evolution or singularity.

\section{The field and its parameters}

Let $(x,y,z)\in\R^3$, $q=x^2+y^2$, and $r=\sqrt q$. Choose
\[
R\geq0,\qquad a>0,\qquad b>0,\qquad A,B\in\R.
\]
Here $R$ sets the ring radius, $a$ controls radial concentration, $b$ controls
vertical extent, $A$ sets the meridional flow strength, and $B$ sets the swirl.
If the coordinates have units of length, $R,a,b$ have those same units and
$A,B$ have units of inverse time. All numerical examples below use
nondimensional variables.

Define
\begin{align}
\phi(r,z)&=\exp\!\left(-\frac{(r^2-R^2)^2}{a^4}-\frac{z^2}{b^2}\right),
\\
\psi(r,z)&=A r^2z\phi(r,z),\qquad u_\theta=B r\phi(r,z).
\end{align}
For $r>0$, the stream-function relations give
\begin{align}
u_r&=-\frac1r\partial_z\psi
     =-A r\left(1-\frac{2z^2}{b^2}\right)\phi,\\
u_z&=\frac1r\partial_r\psi
     =A z\left(2-\frac{4r^2(r^2-R^2)}{a^4}\right)\phi.
\end{align}
These follow from
$\partial_z\phi=-2z\phi/b^2$ and
$\partial_r\phi=-4r(r^2-R^2)\phi/a^4$.

Set
\[
H=1-\frac{2z^2}{b^2},\qquad
C=2-\frac{4q(q-R^2)}{a^4}.
\]
The Cartesian field, which is the definition to use on the axis as well, is
\begin{equation}\label{eq:cartesian}
\boxed{
\begin{aligned}
u_x&=(-A xH-B y)\phi,\\
u_y&=(-A yH+B x)\phi,\\
u_z&=A zC\phi.
\end{aligned}}
\end{equation}
There is no division by $r$ in \eqref{eq:cartesian}. In particular,
\[
\uu(0,0,z)=
\left(0,0,\,2Az\exp\!\left[-\frac{R^4}{a^4}-\frac{z^2}{b^2}\right]\right).
\]
The earlier correction is the special case $a=b=\alpha$ and $A=B=1$.

\section{Smoothness, decay, and incompressibility}

\begin{proposition}
For every allowed parameter choice, $\uu\in\mathcal S(\R^3;\R^3)$ and
$\nabla\cdot\uu=0$ everywhere. Consequently its kinetic energy and the
$L^2$ norm of every spatial derivative are finite.
\end{proposition}

\begin{proof}
The exponent
\[
F(x,y,z)=\frac{(x^2+y^2-R^2)^2}{a^4}+\frac{z^2}{b^2}
\]
is a polynomial. Each component of $\uu$, and each of its derivatives, is a
polynomial times $e^{-F}$. In particular the field is real analytic in
Cartesian coordinates, including on the axis.

To check rapid decay, use
\[
(q-R^2)^2\geq\frac12q^2-R^4.
\]
Thus
\[
0<\phi\leq e^{R^4/a^4}
\exp\!\left(-\frac{q^2}{2a^4}-\frac{z^2}{b^2}\right).
\]
This decay dominates every polynomial in $x,y,z$. For all multi-indices
$\beta,\gamma$ and each component $j$,
\[
\sup_{(x,y,z)\in\R^3}
\big| (x,y,z)^\beta\partial^\gamma u_j(x,y,z)\big|<\infty.
\]
This is the Schwartz-space condition.

For divergence, note
\[
x\partial_x\phi+y\partial_y\phi
=-\frac{4q(q-R^2)}{a^4}\phi,\qquad
-By\partial_x\phi+Bx\partial_y\phi=0.
\]
Differentiating \eqref{eq:cartesian} therefore gives
\begin{align*}
\partial_xu_x+\partial_yu_y
&=-2AH\phi+\frac{4AHq(q-R^2)}{a^4}\phi=-AHC\phi,\\
\partial_zu_z
&=AC(\phi+z\partial_z\phi)=ACH\phi.
\end{align*}
The terms cancel identically. This Cartesian calculation holds at $r=0$
as well as at $r>0$.
\end{proof}

The swirl does not change this cancellation. Its direction is tangent to
circles about the axis and its magnitude has no angular dependence.

\section{Exact kinetic energy}

Use energy per unit constant density:
\[
E_0=\frac12\int_{\R^3}|\uu|^2\dd x\,\dd y\,\dd z.
\]
Multiplying by the physical density gives the physical kinetic energy.
Introduce
\begin{align*}
g(q)&=\exp\!\left(-\frac{2(q-R^2)^2}{a^4}\right),&
Z_0&=\int_{\R}e^{-2z^2/b^2}\dd z=b\sqrt{\frac{\pi}{2}},\\
M_0&=\int_0^\infty g(q)\dd q,&
M_1&=\int_0^\infty qg(q)\dd q.
\end{align*}
The two radial moments are explicit:
\begin{align}
M_0&=\frac{a^2\sqrt\pi}{2\sqrt2}
\left[1+\operatorname{erf}\!\left(\frac{\sqrt2 R^2}{a^2}\right)\right],\\
M_1&=R^2M_0+\frac{a^4}{4}e^{-2R^4/a^4}.
\end{align}
Here $\operatorname{erf}(s)=\frac{2}{\sqrt\pi}\int_0^s e^{-t^2}\dd t$.

\begin{proposition}
The kinetic energy of \eqref{eq:cartesian} is
\begin{equation}\label{eq:energy}
\boxed{
E_0=\frac{\pi Z_0}{2}
\left[
\left(B^2+\frac{3A^2}{4}+\frac{A^2b^2R^2}{a^4}\right)M_1
+\frac{3A^2b^2}{4}M_0
\right].}
\end{equation}
\end{proposition}

\begin{proof}
The squared cylindrical components add to
\[
|\uu|^2=
\left[A^2r^2H^2+B^2r^2+A^2z^2C^2\right]\phi^2.
\]
Use $r\dd r=\dd q/2$ and the Gaussian moments
\[
\int_{\R}H^2e^{-2z^2/b^2}\dd z=\frac34 Z_0,\qquad
\int_{\R}z^2e^{-2z^2/b^2}\dd z=\frac{b^2}{4}Z_0.
\]
The angular integral is $2\pi$, so
\begin{equation}\label{eq:energy-pre}
E_0=\frac{\pi Z_0}{2}
\int_0^\infty
\left[
\left(\frac{3A^2}{4}+B^2\right)q+\frac{A^2b^2}{4}C(q)^2
\right]g(q)\dd q.
\end{equation}

It remains to simplify $K=\int_0^\infty C(q)^2g(q)\dd q$.
Since $g'=-4(q-R^2)g/a^4$ and $C=2+qg'/g$,
\[
K=4M_0+4\int_0^\infty qg'\dd q
  +\int_0^\infty q^2\frac{(g')^2}{g}\dd q.
\]
The first two terms cancel by integration by parts. Also
$(g')^2/g=g''+4g/a^4$. Writing $M_2=\int_0^\infty q^2g(q)\dd q$ yields
\[
K=2M_0+\frac4{a^4}M_2,\qquad
M_2=R^2M_1+\frac{a^4}{4}M_0.
\]
All boundary terms with $q$ or $q^2$ vanish at zero and infinity.
Consequently
\[
K=3M_0+\frac{4R^2}{a^4}M_1.
\]
Substitution in \eqref{eq:energy-pre} proves \eqref{eq:energy}.
\end{proof}

\subsection*{A reproducible example and energy control}

For $R=1.5$, $a=b=1$, and $A=B=1$,
\[
E_0=24.05715786588763.
\]
Scaling both rates by $\lambda$ scales the velocity by $\lambda$ and the
energy by $\lambda^2$. To set a target $E_*>0$ for a nonzero base field, use
\[
\lambda=\sqrt{\frac{E_*}{E_0}},\qquad
(A,B)\longmapsto(\lambda A,\lambda B).
\]
The example has unit energy when
$A=B=0.20388150977663805$.

\section{What changed from the earlier radial Gaussian}

The old envelope was
\[
\phi_{\rm old}(r,z)=
\exp\!\left(-\frac{(r-R)^2+z^2}{\alpha^2}\right).
\]
On $y=0$, its radial coordinate is $r=|x|$. For $R>0$, the one-sided
$x$ derivatives at $x=0$ are
\[
D_x^+\phi_{\rm old}=\frac{2R}{\alpha^2}\phi_{\rm old}(0,z),
\qquad
D_x^-\phi_{\rm old}=-\frac{2R}{\alpha^2}\phi_{\rm old}(0,z).
\]
They do not agree. The old velocity itself also fails to be smooth in
general: for fixed $z$,
\[
u_{z,\rm old}(x,0,z)=
Az\phi_{\rm old}(0,z)
\left(2+\frac{6R}{\alpha^2}|x|+O(x^2)\right).
\]
When $AzR\ne0$, this has a cusp. The new envelope removes the problematic
$|x|$ dependence by using $r^2=x^2+y^2$ inside a polynomial.

The replacement changes the ring shape. For $R>0$ and $r=R+\eta$,
\[
\frac{(r^2-R^2)^2}{a^4}
=\frac{4R^2\eta^2+4R\eta^3+\eta^4}{a^4}.
\]
The local radial Gaussian width, in the convention
$e^{-(\eta/w)^2}$, is approximately $w=a^2/(2R)$.
To match the old width $\alpha$ to quadratic order at the ring, choose
$a=\sqrt{2R\alpha}$ and $b=\alpha$.
This matches local curvature at the peak, not the entire original profile.
For $R=0$, the radial envelope is $e^{-r^4/a^4}$ and is still smooth.

\section{Verification and the remaining scope}

The accompanying \texttt{initial\_field.py} evaluates the Cartesian formula
directly and implements the exact energy and amplitude normalization.
Seven tests cover the axis values, numerical derivatives of the stream
function, Cartesian divergence, rotation about the axis, energy scaling,
the zero field, and invalid geometric parameters.

An independent numerical check integrates $|\uu|^2$ in cylindrical
coordinates with a 180-point Gauss--Legendre rule in each variable.
Four parameter choices, including $R=0$ and unequal widths, agree with
\eqref{eq:energy} to relative error below $1.3\times10^{-14}$ in floating
point. The integration was truncated at
$r^2=R^2+8a^2$ and $|z|=8b$.
These checks validate the implementation; the identities and estimates
above establish the mathematical properties for the full parameter family.

If a scalar initial field is wanted, $S_0=S_*\phi$ is also Schwartz for any
finite constant $S_*$. That fact alone does not ensure smoothness of a
force involving $|\nabla S|$ at critical points of $S$.
The forcing law and the accuracy of the evolution scheme still require
separate analysis.

For the exact initial-data and forcing requirements in the standard
Navier--Stokes problem, see Charles Fefferman's
\href{https://www.claymath.org/library/monographs/MPPc.pdf}
{official problem description for the Clay Mathematics Institute}.

\section{What the stream-function tests establish}

The 20 submitted checks for the original radial Gaussian pass, including
finite-difference divergence away from the axis, reflection in $z$, and
bounded radial velocity near $r=0$. These properties are consistent with the
axis cusp derived above: boundedness alone does not imply differentiability.
For $R=1.5$, $\alpha=1$, $A=1$ and $z=0.4$, the one-sided $x$ slopes of the
old $u_z$ approach opposite nonzero values
$\pm 3.6e^{-2.41}$. The corrected field has matching zero slopes there.

For $r>0$, the contour-tangency identity follows directly:
\[
u_r\partial_r\psi+u_z\partial_z\psi
=-\frac{\partial_z\psi\,\partial_r\psi}{r}
 +\frac{\partial_r\psi\,\partial_z\psi}{r}=0.
\]
Changing the initial swirl keeps these meridional contours fixed, while
changing the azimuthal motion and generally the subsequent evolution.
The corrected envelope is symmetric about $q=R^2$ as a function of
$q=r^2$ (where both comparison points have $q\geq0$); it is generally
not symmetric about $r=R$ as a function of $r$.

\section{A compatible projection on a periodic grid}

Let $h=L/N$ be the grid spacing, let $D_h$ denote centered-difference
divergence, and let $G_h$ denote the matching centered gradient. Given
the velocity $\mathbf u^*$ after the original explicit update, solve
\[
D_hG_h\chi=D_h\mathbf u^*,\qquad
\mathbf u^{n+1}=\mathbf u^*-G_h\chi.
\]
Here $\chi$ is a velocity correction potential. In this nonincremental
time-step construction, $\rho\chi/\Delta t$ is a discrete pressure
approximation. An initial projection has no assigned physical time step.

Because both first derivatives are centered,
\[
(D_hG_h\chi)_{ijk}=
\frac{\chi_{i+2,j,k}+\chi_{i-2,j,k}
      \chi_{i,j+2,k}+\chi_{i,j-2,k}
      \chi_{i,j,k+2}+\chi_{i,j,k-2}
      -6\chi_{ijk}}{4h^2},
\]
with periodic indices. Thus, for $b_{ijk}=(D_h\mathbf u^*)_{ijk}$,
the undamped Jacobi target is
\[
\chi^{\mathrm{target}}_{ijk}
=\frac{\sum_{\text{six neighbors at distance }2h}\chi
       -4h^2b_{ijk}}{6}.
\]
The submitted projection omitted the factor $h^2$. On modes where that
iteration converges, the incorrect solve leaves
$(1-h^{-2})D_h\mathbf u^*$ as the divergence. For $h=0.375$ and $0.1875$,
the residual magnitude can therefore be about $6.11$ and $27.44$ times
the starting value.

There is a second issue: when $N$ is divisible by four, the mode
$\cos[\pi(i+j+k)/2]$ gives an undamped Jacobi error multiplier of $-1$.
It oscillates. The standard-library implementation uses
\[
\chi^{(\ell+1)}
=\tfrac13\chi^{(\ell)}+\tfrac23\chi^{\mathrm{target}}
\]
and checks the maximum residual against an absolute plus relative tolerance.
Failure to converge is reported. The optional Fourier implementation used
for the box experiments solves the same discrete equation with symbol
\[
\lambda(\mathbf m)
=-\frac1{h^2}\sum_{d=1}^3\sin^2(2\pi m_d/N).
\]
It sets the potential to zero on the null modes. For even $N$ there are
eight combinations with each $m_d$ equal to $0$ or $N/2$.
The centered divergence and gradient cannot see these checkerboard modes;
small measured divergence alone does not exclude them.

On the periodic grid $G_h=-D_h^*$ under the grid inner product. An exact
compatible solve therefore gives an orthogonal projection: it preserves
componentwise mean velocity and cannot increase the discrete kinetic
energy during the projection alone. This says nothing by itself about
stability of the preceding explicit update.
The numerical tests check these properties, non-unit spacing, odd and even
grids, and the oscillating mode.

For the historical projection-method construction, see
\href{https://math.berkeley.edu/~chorin/chorin68.pdf}
{A.~J. Chorin, Numerical Solution of the Navier--Stokes Equations (1968)}.
The centered operator and its limitations here are derived explicitly above.

The repository review also found two Copilot variants that solve a
nearest-neighbor pressure Poisson equation and then subtract a centered
gradient. That nearest-neighbor Laplacian is not $D_hG_h$ for the centered
operators used here. Even a converged solve of that different equation need
not remove the measured divergence. These are implementation variants of
the same projection task; they should be reconciled rather than maintained
as independent solvers.

\section{The small-box experiment}

All grids sample cell centers
$x_i=(i-(N-1)/2)L/N$, and similarly for $y,z$. Opposite faces are identified
periodically. Sampling a divergence-free whole-space field need not give
$D_h\mathbf u=0$, and a rapidly decreasing field is not automatically periodic.
No artificial radius clamp is used on the axis.

For the original field, $R=1.5$ and $\alpha=1$, the raw starting diagnostics
on $L=6$ are:
\begin{center}
\begin{tabular}{r r r r r}
\hline
$N$ & $h$ & $\max|D_h\mathbf u|$ &
interior maximum & $\max\|\nabla_h\mathbf u\|_F$\\
\hline
8  & 0.7500 & 0.46351 & 0.46351 & 2.02332\\
16 & 0.3750 & 0.91375 & 0.23399 & 3.39830\\
32 & 0.1875 & 1.71323 & 0.07559 & 3.85268\\
\hline
\end{tabular}
\end{center}
The interior maximum excludes the outermost cell layer, where centered
differences wrap around the boundary. For $N=16,32$ the global divergence
maximum lies on that layer. Interior error decreases while the boundary
error grows. A nonzero jump across a periodic seam produces derivative
terms proportional to the jump divided by $h$.

To separate box size from resolution, compare equal spacing $h=0.375$:
\begin{center}
\begin{tabular}{r r r r}
\hline
$L$ & $N$ & raw $\max|D_h\mathbf u|$ &
sampled opposing-face velocity mismatch\\
\hline
6  & 16 & 0.91375 & $0.81910$\\
18 & 48 & 0.23399 & $8.58\times10^{-24}$\\
\hline
\end{tabular}
\end{center}
Here the mismatch is the largest Euclidean velocity difference at opposite
physical faces sampled at the tangential grid coordinates. Enlarging the
box at fixed $N$ also enlarges $h$ and under-resolves the ring: for $L=18$,
$N=16,32$ the raw divergence maxima are $0.22016,0.40763$ respectively.
Their nonmonotonic behavior is not evidence of resolved convergence.

\subsection*{Four steps and a common smooth start}

The experiment uses the original explicit advection and diffusion update,
$\nu=0.01$, $S=0$, $P_U=0$, and sets the force coefficient $\sigma=0$.
The last setting removes the tiny nonzero tail of the original tanh force.
The original-Gaussian reproduction starts both branches from the raw samples.
After four steps of size $\Delta t=0.02$, the $L=6$ divergence maxima are:
\begin{center}
\begin{tabular}{r r r}
\hline
$N$ & without projection & with corrected projection\\
\hline
8  & 0.48463 & $2.22\times10^{-16}$\\
16 & 0.62311 & $6.66\times10^{-16}$\\
32 & 1.23564 & $1.78\times10^{-15}$\\
\hline
\end{tabular}
\end{center}
The first projected step changes the discrete starting field as well as
correcting divergence created by evolution.

A second comparison uses the smooth field with
$R=1.5$, $a=\sqrt3$, $b=1$, $A=B=1$, matching the original local ring width.
It projects the sampled start once, then copies that same start to both
branches. At $L=6$, $N=32$, this initial correction has relative grid
$L^2$ size $0.009636$. At time $0.08$, the divergence maxima are
$1.05212$ without projection and $1.33\times10^{-15}$ with projection.
The maximum gradient norms are $4.05142$ and $4.08213$ respectively.
Projection enforces the chosen discrete divergence; it need not lower the
maximum gradient.

Halving the time step to $0.01$ and taking eight steps to the same time
gives $1.03447$ and $1.33\times10^{-15}$ for the two divergence maxima.
The projected maximum gradient changes from $4.08213$ to $4.07891$,
and its discrete energy from $38.80037$ to $38.77337$.
This is one time-resolution check, not a convergence or stability proof.

Even the smooth field has a small nonzero periodic face mismatch on $L=6$.
The next spatial-accuracy experiment should make the boundary treatment
explicit (for example, periodize the rapidly decreasing field), then refine
space and time with that choice fixed. A finite grid and four time steps
cannot establish continuum regularity or singularity.

\end{document}
